\documentclass[conference]{IEEEtran}

\usepackage{amsmath}
\usepackage{booktabs}
\usepackage{cite}
\usepackage{graphicx}
\usepackage{url}

\graphicspath{{../figures/submission/}}

\begin{document}

\title{Contact-Rich Diff-MPPI: Hybrid Sampling and Gradient Refinement
for GPU Manipulation}

\author{\IEEEauthorblockN{Anonymous Authors}}

\maketitle

\begin{abstract}
Sampling-based model predictive control searches broadly but may spend many
rollouts resolving a local contact mode.  Differentiable trajectory
optimization supplies local information, yet can fail when initialized in the
wrong basin.  We study a minimal, training-free hybrid that preserves the
standard Model Predictive Path Integral (MPPI) update and then applies one or
three clipped autodiff steps to the nominal control sequence.  The evaluation
keeps failures visible and separates three claims.  First, a preregistered
32,400-episode GPU matrix spans 12 model and geometry conditions, five
contact-rich tasks, six planners, three rollout budgets, and 30 paired seeds;
Diff-MPPI-3 produces 33 Holm-significant positive and 6 negative cells versus
MPPI.  Second, in enforced 10 ms control slots with disjoint calibration and
held-out seeds, Diff-MPPI-3 reaches 1.00 real-time success on contact loss
versus 0.467 for MPPI, a paired difference of \(+0.533\) with Holm
\(p=0.000305\).  Third, 3,150 closed-loop MuJoCo episodes test an independently
executed contact plant across friction, mass, and observation-noise changes;
Diff-MPPI-3 reaches 0.457 aggregate success versus 0.289 for MPPI, with three
positive and zero negative Holm-significant cells.  Tall-object regressions
and an unsolved detour family bound the result.  The evidence supports a
localized compute--quality advantage for selected contact-rich tasks, not
universal planner dominance or real-robot transfer.
\end{abstract}

\begin{IEEEkeywords}
model predictive control, MPPI, differentiable control, contact-rich
manipulation, GPU computing
\end{IEEEkeywords}

\section{Introduction}

Model Predictive Path Integral control uses importance-weighted stochastic
rollouts to optimize nonlinear control sequences and maps naturally to GPU
execution~\cite{williams2017mppi}.  Its broad search is valuable for
nonconvex costs and contact-mode changes, but finite samples can repeatedly
miss the narrow local corrections needed to restore or maintain contact.
Gradient-based refinement has complementary behavior: once a nominal sequence
enters an informative basin, local derivatives can sharpen it cheaply, while
a wrong basin or a topological detour can still defeat the method.

We investigate the smallest useful combination of these mechanisms.
Diff-MPPI leaves the sampling update intact and appends a short, clipped
forward-mode autodiff refinement.  It learns no sampler, uses no offline
dataset, and does not differentiate through the sampling distribution.  The
question is deliberately narrow: does this local refinement improve selected
contact-rich tasks under fixed samples, enforced deadlines, and an
independently executed contact plant?

The evaluation is organized around paired, content-addressed evidence instead
of a collection of favorable examples.  It contributes:

\begin{enumerate}
  \item a minimal GPU hybrid of standard MPPI sampling and one or three local
  autodiff steps for planar contact-rich manipulation;
  \item a 32,400-episode robustness protocol with paired seeds, full-family
  Holm correction, smooth and hard-contact plants, and retained negative
  cells;
  \item an enforced 10 ms matched-compute evaluation with disjoint calibration
  and held-out seeds; and
  \item a 3,150-episode closed-loop MuJoCo transfer study in which the planners
  retain their nominal smooth model while a distinct engine advances the
  executed plant.
\end{enumerate}

The result is not universal.  Diff-MPPI-3 has six corrected regressions under
a tall-box condition, and every method fails the detour family in the
deadline-matched experiment.  These failures are part of the declared result.

\section{Related Work}

\textbf{Sampling-based MPC.}
MPPI updates a nominal control sequence by exponentially weighting sampled
trajectory perturbations~\cite{williams2017mppi}.  The update tolerates
nonsmooth costs and parallelizes over rollouts, which makes it a strong
baseline for real-time GPU control.

\textbf{Sampling and local optimization.}
CEM-GD combines population search with gradient-based refinement for
model-based planning~\cite{bharadhwaj2020cemgd}.  Other hybrids add feedback,
learned proposals, or trajectory-optimization subproblems.  Our contribution
is not the generic idea of combining sampling and gradients.  It is the
controlled localization of a minimal post-MPPI refinement under a declared
contact-family statistical protocol and a wall-clock control slot.

\textbf{Differentiable contact and external execution.}
Contact gradients can expose local information that random perturbations
rarely sample, but the gradient depends on a model whose smoothing and contact
law may be wrong.  We therefore use a smooth nominal model for optimization
and separately execute commands in MuJoCo~\cite{todorov2012mujoco}.  This is
closed-loop sim-to-sim transfer rather than open-loop replay, a standard robot
benchmark, or physical-robot evidence.

\section{Method}

\subsection{MPPI update}

Let \(U=(u_0,\ldots,u_{T-1})\) be the nominal sequence and
\(\epsilon_k\) a sampled perturbation.  A rollout produces cost \(S_k\).
With temperature \(\lambda\) and \(\rho=\min_k S_k\), normalized weights are
\begin{equation}
w_k =
\frac{\exp(-(S_k-\rho)/\lambda)}
     {\sum_j \exp(-(S_j-\rho)/\lambda)} ,
\end{equation}
and the standard update is
\begin{equation}
U^+ = U + \sum_k w_k\epsilon_k .
\end{equation}
Rollouts, stage costs, and weighted reductions execute on the GPU.

\subsection{Local autodiff refinement}

Starting from \(U^{(0)}=U^+\), the hybrid differentiates the nominal smooth
contact rollout and applies
\begin{equation}
U^{(j+1)} =
\Pi_{\mathcal U}\!\left[
U^{(j)}-\alpha\,
\operatorname{clip}\!\left(\nabla_U J(U^{(j)}),g_{\max}\right)
\right],
\end{equation}
where \(\Pi_{\mathcal U}\) enforces control limits.  Diff-MPPI-1 takes one step
and Diff-MPPI-3 takes three.  The controller executes the first refined
command, shifts the horizon, and repeats from the observed next state.

The derivative is local to the post-MPPI nominal sequence.  We do not
differentiate through random-number generation, sampled rollouts, importance
weights, or the sampling update.  Sampling remains responsible for mode-level
search; refinement only sharpens the selected mode.

\subsection{Contact models and baselines}

The nominal planner uses a differentiable planar box-contact approximation.
The robustness matrix varies contact gain, object size and aspect ratio, and
also includes a momentum-carrying hard-contact plant with friction and damping
sweeps.  The broad comparison contains vanilla MPPI, Diff-MPPI-1,
Diff-MPPI-3, SOPPI, SOPPI-fast, and a hard-model MPPI reference.  SOPPI-fast
contains one nominal gradient step and is therefore not labeled a pure
sampling-only or pure-SVGD baseline.

\section{Experimental Protocol}

\subsection{Tasks and broad robustness matrix}

Five planar box-manipulation tasks cover sustained swivel, strict pose
alignment, obstacle detour, recovery after contact loss, and curved contact.
The release matrix comprises 12 model/geometry conditions, five tasks, six
planners, \(K\in\{128,256,512\}\), and 30 paired seeds per cell: 32,400
episodes and 1,080 summary cells.

Comparisons use paired success differences, paired bootstrap 95\% confidence
intervals, exact McNemar tests, and Holm correction over the declared family.
Statistical significance is not an artifact-validity condition; losses and
zero-success cells remain valid evidence.

\subsection{Enforced 10 ms control slot}

Calibration and evaluation seeds are disjoint.  For three planners and five
tasks, 25 calibration seeds select the largest sample count with zero
calibration deadline misses, giving 375 calibration episodes.  Thirty held-out
seeds per planner/task cell give 450 evaluation episodes.  Every planner
receives an enforced 10 ms slot, and real-time success requires both task
success and zero deadline misses.  All three planners select \(K=1024\) on the
reference GPU; fairness comes from the enforced slot, not unequal \(K\).

\subsection{Closed-loop MuJoCo transfer}

The external-fidelity block contains 3,150 closed-loop episodes over seven
nominal, friction, mass, and observation-noise conditions; five tasks; three
planners; \(K=256\); and 30 paired seeds.  CUDA planners optimize the nominal
smooth dynamics.  MuJoCo 3.11.0 advances every chosen command and returns the
next state.  The same paired bootstrap, McNemar, and full-family Holm procedure
is applied.

\subsection{Evidence freeze}

All frozen results were collected on one NVIDIA GeForce GTX 1660 Ti.  Each
experiment block retains its source commit, hardware identity, commands,
matrix shape, raw artifact hashes, report hashes, and independent validator
result.  The supplementary ledger reruns numerical assertions against the
copied evidence.  Independent hardware replication is desirable but is not
implied.

\section{Results}

\subsection{Broad robustness}

\begin{table}[t]
\centering
\caption{Aggregate 32,400-episode robustness matrix.}
\label{tab:robustness}
\begin{tabular}{lrrr}
\toprule
Planner & Episodes & Success & Mean ms \\
\midrule
Diff-MPPI-3    & 5,400 & 0.614 & 2.655 \\
Diff-MPPI-1    & 5,400 & 0.595 & 0.959 \\
SOPPI-fast     & 5,400 & 0.557 & 1.464 \\
MPPI           & 5,400 & 0.453 & 0.120 \\
SOPPI          & 5,400 & 0.446 & 0.426 \\
MPPI-hardmodel & 5,400 & 0.331 & 0.159 \\
\bottomrule
\end{tabular}
\end{table}

Table~\ref{tab:robustness} shows the aggregate ordering.  Across the 360 paired
comparisons against MPPI, the declared family contains 33 Holm-significant
positive and 6 negative success cells.  All six negative cells are
Diff-MPPI-3 under the tall-box condition.  The hard-contact friction/damping
conditions retain positive cells under a structurally different,
momentum-carrying plant.

\begin{figure}[t]
\centering
\includegraphics[width=\columnwidth]{contact_robustness.pdf}
\caption{Corrected positive and negative Diff-MPPI-3 cells versus MPPI across
the robustness conditions.}
\label{fig:robustness}
\end{figure}

\subsection{Deadline-matched control}

\begin{table}[t]
\centering
\caption{Held-out episodes under the enforced 10 ms slot.}
\label{tab:deadline}
\begin{tabular}{lrrr}
\toprule
Planner & Episodes & RT success & Misses \\
\midrule
Diff-MPPI-3 & 150 & 0.800 & 1 \\
SOPPI-fast  & 150 & 0.793 & 55 \\
MPPI        & 150 & 0.673 & 4 \\
\bottomrule
\end{tabular}
\end{table}

The aggregate row mixes tasks of different difficulty, so the preregistered
paired family is primary.  On contact loss, Diff-MPPI-3 reaches 1.00
real-time success versus 0.467 for MPPI.  The paired difference is \(+0.533\),
with bootstrap 95\% CI \([+0.367,+0.700]\) and Holm
\(p=0.000305\).  SOPPI-fast reaches 0.967 and improves by \(+0.500\), with
Holm \(p=0.002472\).  The other four task cells are not Holm-significant, and
all planners remain 0/30 on detour.

\begin{figure}[t]
\centering
\includegraphics[width=\columnwidth]{contact_matched_compute.pdf}
\caption{Held-out real-time success under identical 10 ms control slots.}
\label{fig:matched}
\end{figure}

\subsection{Closed-loop MuJoCo transfer}

\begin{table}[t]
\centering
\caption{Aggregate closed-loop MuJoCo transfer.}
\label{tab:mujoco}
\begin{tabular}{lrrr}
\toprule
Planner & Episodes & Success & Mean ms \\
\midrule
Diff-MPPI-3 & 1,050 & 0.457 & 2.614 \\
SOPPI-fast  & 1,050 & 0.340 & 1.432 \\
MPPI        & 1,050 & 0.289 & 0.117 \\
\bottomrule
\end{tabular}
\end{table}

The 70-cell paired family contains three Holm-significant positive
Diff-MPPI-3 cells and zero negative cells.  No individual
observation-noise cell survives full-family correction, so sensing-noise
effects remain descriptive.  Detour remains zero-success across planners and
conditions.

\begin{figure}[t]
\centering
\includegraphics[width=\columnwidth]{contact_external_fidelity.pdf}
\caption{Closed-loop MuJoCo success across physical and sensing variations.}
\label{fig:mujoco}
\end{figure}

\section{Discussion and Limitations}

The results support a division of labor.  Sampling performs broad search.
When the nominal sequence enters a useful contact basin, one or three gradient
steps can correct alignment or contact loss without replacing MPPI.  The
contact-loss result survives an enforced control slot, and the aggregate
advantage transfers to an independently executed MuJoCo plant.

The failures define the boundary.  Tall objects produce six corrected
negative cells, which indicates sensitivity to geometry or gradient scaling.
Detour remains unsolved because obstacle-induced topology requires a different
mode rather than sharper local refinement.  SOPPI-fast is competitive, but
its own gradient step prevents a pure sampling-versus-gradient interpretation.

All frozen experiments use one GTX 1660 Ti.  The nominal and hard-contact
plants are custom GPU simulations.  The external plant is a custom planar
MuJoCo model, not a standard manipulator suite.  Closed-loop MuJoCo transfer
is sim-to-sim evidence, not real-robot evidence.  The work does not cover
deformable contact, 3D grasping, or physical hardware safety constraints.

\section{Conclusion}

A short, training-free autodiff refinement can improve selected contact-rich
MPPI cells under fixed samples, enforced deadlines, and sim-to-sim model
mismatch.  The 32,400-episode robustness family, exact 10 ms evaluation, and
3,150-episode closed-loop MuJoCo block also expose corrected regressions and
an unsolved detour family.  The supported claim is therefore a localized
compute--quality advantage, not universal planner dominance.

\bibliographystyle{IEEEtran}
\bibliography{references}

\end{document}
